\section{Resultados y límites de Search-o1, y Search-R1}

\subsection{El desempeño sigue escalando conforme aumentan los documentos}

\videofigure{figures/search-o1-scaling.jpg}{Search-o1 mejora de manera continua en múltiples dominios científicos conforme aumenta el número de documentos recuperados, mientras que el RAG estándar se satura rápidamente.}{01:00:34--01:01:28}

\videofigure{figures/gpqa-results.jpg}{En el conjunto extendido de GPQA, Search-o1 supera al razonamiento directo y al RAG estándar, y reduce la brecha con expertos humanos.}{01:01:28--01:03:18}

\videofigure{figures/retrieved-doc-scaling.jpg}{Al pasar de top-1 a top-10 documentos, Search-o1 puede aprovechar más evidencia; el RAG común se ve arrastrado más fácilmente por el ruido.}{01:03:18--01:04:38}

\videofigure{figures/multihop-qa.jpg}{En preguntas y respuestas de varios saltos como HotpotQA, 2Wiki, MuSiQue y Bamboogle, la búsqueda iterativa supera a la recuperación de un solo salto.}{01:04:38--01:06:12}

Estos resultados respaldan un principio de diseño de sistemas: la escala de la recuperación debe crecer junto con la capacidad de depurar evidencia. Aumentar solo el número de documentos sin mejorar el query refinement y la compression no produce ganancias estables.

\subsection{Reducir la incertidumbre no equivale a que los hechos sean del todo confiables}

\videofigure{figures/search-o1-insights.jpg}{El artículo analiza el origen de las ganancias de Search-o1 desde tres ángulos: palabras de incertidumbre, escala de documentos y consultas de múltiples rondas.}{01:06:12--01:06:46}

\videofigure{figures/reasoning-quality.jpg}{Al incorporar la búsqueda, el lenguaje ambiguo en el razonamiento disminuye y aumentan los enunciados respaldados por evidencia.}{01:06:46--01:08:02}

\begin{warningbox}{«Sonar más seguro» no es condición suficiente de corrección}
El retriever puede devolver páginas web erróneas, el compresor puede malinterpretar la evidencia, y el modelo también puede vincular incorrectamente una cita con una conclusión. Sigue siendo necesaria la calidad de las fuentes, la claim-level citation y la verificación de la respuesta final.
\end{warningbox}

\subsection{Search-o1 y Search-R1}

Search-o1 enseña principalmente al modelo cuándo buscar mediante prompting y un agent loop explícito; Search-R1, en cambio, busca aprender la estrategia de búsqueda dentro de la policy usando reinforcement learning. Las diferencias entre ambos pueden resumirse así:

\begin{table}[H]
\centering
\small
\begin{tabular}{>{\raggedright\arraybackslash}p{3cm}>{\raggedright\arraybackslash}p{4.7cm}>{\raggedright\arraybackslash}p{4.7cm}}
\toprule
\textbf{Dimensión} & \textbf{Search-o1} & \textbf{Search-R1} \\
\midrule
Forma de control & prompt, tokens especiales y un bucle escrito a mano & RL que aprende cuándo buscar y cómo usar los resultados \\
Ventajas & fácil de depurar, se puede sustituir el componente de recuperación, no requiere entrenamiento & puede reducir reglas manuales, la policy puede ser más adaptativa \\
Riesgos & el prompt es frágil, la estrategia de invocación no siempre es óptima & el diseño del reward es complejo, puede aprender a buscar de forma oportunista \\
Etapa adecuada & construcción y análisis rápidos de sistemas & cuando se cuenta con un reward confiable y suficiente presupuesto de interacción en línea \\
\bottomrule
\end{tabular}
\end{table}

\begin{importantbox}{El comportamiento de búsqueda en sí también necesita un reward}
Si solo se premia la respuesta final, el agent puede buscar en exceso, no buscar en absoluto o explotar fallas del benchmark. El objetivo ideal también debería considerar la calidad de la respuesta, la cobertura de evidencia, el número de consultas, la latencia y el costo de los documentos.
\end{importantbox}

\subsection{Resumen de la sección}

Search-o1 demuestra que la recuperación de múltiples rondas, bajo demanda y comprimible puede mejorar de manera significativa el razonamiento intensivo en conocimiento; Search-R1 va más allá al tratar la estrategia de invocación como una conducta aprendible. Ambos exigen verificar la evidencia recuperada, en lugar de tratar el retrieval como una verdad natural.

